% ============================================================
% 测试卷 — 第9章 统计
% ============================================================
\documentclass[12pt, a4paper]{ctexart}
\usepackage[margin=2cm]{geometry}
\usepackage{amsmath, amssymb}
\usepackage{enumitem}
\usepackage{fancyhdr}
\usepackage{xcolor}
\usepackage{tikz}
\usepackage{graphicx}
\usepackage{tabularx}
\usepackage{multicol}

\definecolor{mainblue}{RGB}{41, 98, 168}

\pagestyle{fancy}
\fancyhf{}
\fancyhead[L]{\small\textcolor{mainblue}{2025-2026学年度·第9章测试卷}}
\fancyhead[R]{\small\textcolor{mainblue}{数学（必修第二册）}}
\fancyfoot[C]{\small —\ \thepage\ —}
\renewcommand{\headrulewidth}{0.6pt}
\renewcommand{\headrule}{\hbox to\headwidth{\color{mainblue}\leaders\hrule height \headrulewidth\hfill}}

\setCJKmainfont{Hiragino Sans GB}
\setCJKsansfont{Heiti SC}

\newcommand{\blank}[1]{\underline{\hspace{#1}}}

\begin{document}

\begin{center}
{\Large\bfseries 第9章\quad 统计\quad 测试卷}\\[6pt]
{\normalsize 数学（必修第二册）}\\[4pt]
{\small 考试时间：120分钟\quad 满分：150分}
\end{center}

\vspace{2pt}
\begin{center}
\begin{tabular}{|c|c|c|c|c|}
\hline
题号 & 一 & 二 & 三 & 四 & 总分 \\ \hline
得分 & & & & & \\ \hline
\end{tabular}
\end{center}
\vspace{4pt}

\noindent\textbf{注意事项：}
\begin{enumerate}[leftmargin=2em]
  \item 答题前，请在答题卡上填写姓名、班级和考号。
  \item 选择题请用2B铅笔在答题卡上涂黑对应选项。
  \item 非选择题请用黑色签字笔在答题卡指定区域作答。
\end{enumerate}

\vspace{6pt}

% ============ 一、单项选择题 ============
\noindent\rule{\textwidth}{0.4pt}\\[2pt]
{\large\bfseries 一、单项选择题}（本大题共8小题，每小题5分，共40分。每小题只有一个选项符合题目要求。）

\vspace{6pt}

\begin{enumerate}
\item 数据 $2,3,5,7,8$ 的平均数为（\quad）
  \begin{enumerate}[label=\Alph*., itemsep=0pt]
    \item $4$ \quad \item $5$ \quad \item $6$ \quad \item $7$
  \end{enumerate}

\item 数据 $3,1,4,1,5,9,2,6$ 的中位数为（\quad）
  \begin{enumerate}[label=\Alph*., itemsep=0pt]
    \item $3$ \quad \item $3.5$ \quad \item $4$ \quad \item $4.5$
  \end{enumerate}

\item 已知一组数据的平均数为 $10$，方差为 $4$。若将每个数据都扩大为原来的 $2$ 倍，则新数据的方差为（\quad）
  \begin{enumerate}[label=\Alph*., itemsep=0pt]
    \item $4$ \quad \item $8$ \quad \item $12$ \quad \item $16$
  \end{enumerate}

\item 甲组有 $5$ 名同学，平均成绩为 $80$ 分；乙组有 $4$ 名同学，平均成绩为 $85$ 分。两组同学的平均成绩为（\quad）
  \begin{enumerate}[label=\Alph*., itemsep=0pt]
    \item $82$ \quad \item $82.2$ \quad \item $82.5$ \quad \item $83$
  \end{enumerate}

\item 反映一组数据波动大小的统计量是（\quad）
  \begin{enumerate}[label=\Alph*., itemsep=0pt]
    \item 平均数 \quad \item 中位数 \quad \item 方差 \quad \item 众数
  \end{enumerate}

\item 已知一组数据 $x_1,x_2,\ldots,x_n$ 的方差为 $4$，则数据 $2x_1+1,2x_2+1,\ldots,2x_n+1$ 的方差为（\quad）
  \begin{enumerate}[label=\Alph*., itemsep=0pt]
    \item $4$ \quad \item $8$ \quad \item $9$ \quad \item $16$
  \end{enumerate}

\item 某班 $20$ 名同学的数学成绩（满分 $100$ 分）的茎叶图如下：\\
\begin{center}
\begin{tabular}{c|c}
$6$ & $2\ 5\ 8$ \\
$7$ & $1\ 2\ 2\ 3\ 5\ 6\ 8$ \\
$8$ & $2\ 3\ 5\ 5\ 8\ 9$ \\
$9$ & $1\ 4\ 5$ \\
\end{tabular}
\end{center}
\noindent 则下列说法正确的是（\quad）
  \begin{enumerate}[label=\Alph*., itemsep=0pt]
    \item 众数为 $85$
    \item 中位数为 $78.5$
    \item 成绩在 $[80,90)$ 的有 $6$ 人
    \item 平均数大于 $80$
  \end{enumerate}

\item 某校高一年级有 $800$ 名学生，用分层抽样从三个班中抽取 $40$ 名学生。已知一班有 $30$ 人被抽到，一班共有 $600$ 名学生，则高一年级共有的班级数为（\quad）
  \begin{enumerate}[label=\Alph*., itemsep=0pt]
    \item $6$ \quad \item $8$ \quad \item $10$ \quad \item $12$
  \end{enumerate}
\end{enumerate}

\vspace{6pt}

% ============ 二、多项选择题 ============
\noindent\rule{\textwidth}{0.4pt}\\[2pt]
{\large\bfseries 二、多项选择题}（本大题共4小题，每小题5分，共20分。全部选对得5分，部分选对得2分，有选错的得0分。）

\vspace{6pt}

\begin{enumerate}[resume]
\item 下列关于回归分析的说法正确的是（\quad）
  \begin{enumerate}[label=\Alph*., itemsep=0pt]
    \item 回归直线一定过样本点的中心 $(\bar{x},\bar{y})$
    \item 相关系数 $|r|$ 越接近 $1$，线性相关性越强
    \item 回归直线的斜率 $\hat{b}$ 可以为负数
    \item 用回归方程进行预测时，自变量的取值可以任意
  \end{enumerate}

\item 已知一组数据 $x_1,x_2,\ldots,x_{10}$ 的平均数为 $\bar{x}$，方差为 $s^2$，则（\quad）
  \begin{enumerate}[label=\Alph*., itemsep=0pt]
    \item 数据 $x_1+a,x_2+a,\ldots,x_{10}+a$ 的平均数为 $\bar{x}+a$
    \item 数据 $x_1+a,x_2+a,\ldots,x_{10}+a$ 的方差为 $s^2$
    \item 数据 $ax_1,ax_2,\ldots,ax_{10}$ 的平均数为 $a\bar{x}$
    \item 数据 $ax_1,ax_2,\ldots,ax_{10}$ 的方差为 $a^2s^2$
  \end{enumerate}

\item 下列抽样方法中，属于简单随机抽样的是（\quad）
  \begin{enumerate}[label=\Alph*., itemsep=0pt]
    \item 从 $100$ 名同学中随机抽取 $10$ 名参加座谈会
    \item 从全班同学中按学号每隔 $5$ 人抽取一人
    \item 从一副扑克牌中随机抽取 $5$ 张
    \item 从生产线每隔 $10$ 分钟抽取一件产品进行检验
  \end{enumerate}

\item 某次考试成绩服从正态分布 $N(75,100)$，则（\quad）
  \begin{enumerate}[label=\Alph*., itemsep=0pt]
    \item 平均分为 $75$
    \item 标准差为 $10$
    \item 约 $68\%$ 的学生成绩在 $(65,85)$ 之间
    \item 约 $95\%$ 的学生成绩在 $(55,95)$ 之间
  \end{enumerate}
\end{enumerate}

\vspace{6pt}

% ============ 三、填空题 ============
\noindent\rule{\textwidth}{0.4pt}\\[2pt]
{\large\bfseries 三、填空题}（本大题共4小题，每小题5分，共20分。）

\vspace{6pt}

\begin{enumerate}[resume]
\item 已知 $5$ 个数据 $3,5,7,9,x$ 的平均数为 $10$，则 $x=$ \underline{\hspace{2cm}}。

\item 数据 $1,2,3,4,5$ 的方差为 \underline{\hspace{2cm}}，标准差为 \underline{\hspace{2cm}}。

\item 数据 $2,3,a,5,7$ 的中位数为 $4$，则 $a=$ \underline{\hspace{2cm}}。

\item 已知 $x$ 与 $y$ 的回归方程为 $\hat{y}=2x+1$，且 $\bar{x}=3$，则 $\bar{y}=$ \underline{\hspace{2cm}}。
\end{enumerate}

\vspace{6pt}

% ============ 四、解答题 ============
\noindent\rule{\textwidth}{0.4pt}\\[2pt]
{\large\bfseries 四、解答题}（本大题共6小题，共70分。解答应写出文字说明、证明过程或演算步骤。）

\vspace{6pt}

\begin{enumerate}[resume]
\item[17.]（本小题10分）\\
某次数学测试，$10$ 名同学的成绩如下：\\
$62,\ 78,\ 85,\ 92,\ 68,\ 75,\ 88,\ 95,\ 70,\ 82$
\begin{enumerate}[(1)]
  \item 求这组数据的平均数和中位数；
  \item 求这组数据的方差。
\end{enumerate}

\vspace{6cm}

\item[18.]（本小题12分）\\
某班级 $15$ 名同学的数学成绩（满分 $100$ 分）如下：\\
$82,\ 90,\ 75,\ 88,\ 92,\ 65,\ 78,\ 85,\ 96,\ 70,\ 83,\ 91,\ 68,\ 77,\ 86$
\begin{enumerate}[(1)]
  \item 将成绩按区间 $[60,70)$、$[70,80)$、$[80,90)$、$[90,100]$ 进行分组，列出频率分布表；
  \item 求成绩在 $80$ 分以上（含 $80$ 分）的频率。
\end{enumerate}

\vspace{6cm}

\item[19.]（本小题12分）\\
甲、乙两人各射靶 $5$ 次，成绩如下（单位：环）：\\
甲：$8,\ 9,\ 7,\ 10,\ 6$\\
乙：$7,\ 8,\ 8,\ 9,\ 8$
\begin{enumerate}[(1)]
  \item 分别求甲、乙成绩的平均数和方差；
  \item 从平均数和方差两个角度，分析甲、乙两人射击水平的差异。
\end{enumerate}

\vspace{6cm}

\item[20.]（本小题12分）\\
已知 $5$ 组数据 $(x_i,y_i)$ 的统计量如下：$\bar{x}=4$，$\bar{y}=5$，$\displaystyle\sum_{i=1}^{5}x_iy_i=115$，$\displaystyle\sum_{i=1}^{5}x_i^2=90$。
\begin{enumerate}[(1)]
  \item 求回归直线 $\hat{y}=\hat{b}x+\hat{a}$ 的回归系数 $\hat{b}$ 和 $\hat{a}$；
  \item 当 $x=6$ 时，预测 $y$ 的值。
\end{enumerate}

\vspace{6cm}

\item[21.]（本小题12分）\\
从全校 $1200$ 名学生中用分层抽样的方法抽取一个容量为 $60$ 的样本。已知高一年级有 $480$ 人，高二年级有 $360$ 人，高三年级有 $360$ 人。
\begin{enumerate}[(1)]
  \item 求各年级应抽取的人数；
  \item 若高一年级样本中有 $8$ 人成绩优秀，估计高一年级成绩优秀的人数。
\end{enumerate}

\vspace{6cm}

\item[22.]（本小题12分）\\
已知 $x$ 与 $y$ 的 $5$ 组观测数据如下：
\begin{center}
\begin{tabular}{|c|c|c|c|c|c|}
\hline
$x$ & $2$ & $3$ & $4$ & $5$ & $6$ \\ \hline
$y$ & $2.5$ & $3.5$ & $4.2$ & $5.0$ & $5.8$ \\ \hline
\end{tabular}
\end{center}
\begin{enumerate}[(1)]
  \item 求 $\bar{x}$、$\bar{y}$；
  \item 求回归直线方程 $\hat{y}=\hat{b}x+\hat{a}$；
  \item 求相关系数 $r$，并判断 $x$ 与 $y$ 的线性相关程度（$|r|>0.75$ 视为强相关）。
\end{enumerate}

\vspace{6cm}
\end{enumerate}

\vfill
\begin{center}
\small\textcolor{gray}{本试卷仅供教学参考}
\end{center}

\end{document}
